\documentclass[12pt,a4paper]{article}

% ── Packages ──────────────────────────────────────────────────────────────
\usepackage[utf8]{inputenc}
\usepackage[T5]{fontenc}
\usepackage[vietnamese]{babel}
\usepackage{geometry}
\geometry{top=2.5cm, bottom=2.5cm, left=3cm, right=2.5cm}

\usepackage{amsmath, amssymb, amsthm}
\usepackage{mathtools}
\usepackage{algorithm}
\usepackage{algpseudocode}
\usepackage{listings}
\usepackage{xcolor}
\usepackage{graphicx}
\usepackage{float}
\usepackage{booktabs}
\usepackage{array}
\usepackage{tabularx}
\usepackage{multirow}
\usepackage{hyperref}
\usepackage{enumitem}
\usepackage{fancyhdr}
\usepackage{titlesec}
\usepackage{tcolorbox}
\usepackage{tikz}
\usepackage{pgfplots}
\pgfplotsset{compat=1.18}
\usetikzlibrary{arrows.meta, positioning, shapes.geometric, calc}

% ── Màu sắc ───────────────────────────────────────────────────────────────
\definecolor{navyblue}{RGB}{26,35,126}
\definecolor{codegreen}{RGB}{46,125,50}
\definecolor{codegray}{RGB}{128,128,128}
\definecolor{codepurple}{RGB}{106,27,154}
\definecolor{backcolor}{RGB}{245,245,250}
\definecolor{lightblue}{RGB}{232,234,246}

% ── Listing style ─────────────────────────────────────────────────────────
\lstdefinestyle{pythonstyle}{
    backgroundcolor=\color{backcolor},
    commentstyle=\color{codegray}\itshape,
    keywordstyle=\color{navyblue}\bfseries,
    stringstyle=\color{codegreen},
    numberstyle=\tiny\color{codegray},
    basicstyle=\ttfamily\footnotesize,
    breakatwhitespace=false,
    breaklines=true,
    captionpos=b,
    keepspaces=true,
    numbers=left,
    numbersep=5pt,
    showspaces=false,
    showstringspaces=false,
    showtabs=false,
    tabsize=4,
    frame=single,
    rulecolor=\color{navyblue!30},
    language=Python
}
\lstset{style=pythonstyle}

% ── Hyperref setup ────────────────────────────────────────────────────────
\hypersetup{
    colorlinks=true,
    linkcolor=navyblue,
    urlcolor=navyblue,
    citecolor=navyblue,
    pdftitle={Báo cáo Bài tập lớn: Thuật toán Fourier-Motzkin},
    pdfauthor={Sinh viên}
}

% ── Header/Footer ─────────────────────────────────────────────────────────
\pagestyle{fancy}
\fancyhf{}
\fancyhead[L]{\small\textcolor{navyblue}{\textbf{CS112 -- Phân tích và Thiết kế Thuật toán}}}
\fancyhead[R]{\small\textcolor{navyblue}{Fourier-Motzkin LP Solver}}
\fancyfoot[C]{\thepage}
\renewcommand{\headrulewidth}{0.4pt}

% ── Section format ────────────────────────────────────────────────────────
\titleformat{\section}{\large\bfseries\color{navyblue}}{{\thesection.}}{0.5em}{}[\titlerule]
\titleformat{\subsection}{\normalsize\bfseries\color{navyblue!80}}{{\thesubsection.}}{0.5em}{}
\titleformat{\subsubsection}{\normalsize\bfseries\color{navyblue!60}}{{\thesubsubsection.}}{0.5em}{}

% ── Theorem environments ──────────────────────────────────────────────────
\theoremstyle{definition}
\newtheorem{definition}{Định nghĩa}[section]
\newtheorem{theorem}{Định lý}[section]
\newtheorem{lemma}{Bổ đề}[section]
\newtheorem{example}{Ví dụ}[section]
\newtheorem{remark}{Nhận xét}[section]

% ── tcolorbox styles ──────────────────────────────────────────────────────
\tcbuselibrary{skins, breakable}
\newtcolorbox{infobox}[1][]{
    colback=lightblue, colframe=navyblue, fonttitle=\bfseries,
    title=#1, breakable
}
\newtcolorbox{resultbox}[1][]{
    colback=green!5, colframe=codegreen, fonttitle=\bfseries,
    title=#1, breakable
}

% ── Macros ────────────────────────────────────────────────────────────────
\newcommand{\R}{\mathbb{R}}
\newcommand{\FM}{\textsc{Fourier-Motzkin}}
\newcommand{\vx}{\mathbf{x}}
\newcommand{\vc}{\mathbf{c}}
\newcommand{\vb}{\mathbf{b}}
\newcommand{\mA}{\mathbf{A}}

% ══════════════════════════════════════════════════════════════════════════
\begin{document}
% ══════════════════════════════════════════════════════════════════════════

% ── Trang bìa ─────────────────────────────────────────────────────────────
\begin{titlepage}
\centering
\vspace*{1cm}

{\Large\bfseries TRƯỜNG ĐẠI HỌC KHOA HỌC TỰ NHIÊN\\
ĐẠI HỌC QUỐC GIA TP. HỒ CHÍ MINH}

\vspace{0.5cm}
{\large Khoa Công nghệ Thông tin}

\vspace{1.5cm}
\includegraphics[width=3cm]{logo.png}  % Thay bằng logo trường nếu có
\vspace{1cm}

\begin{tcolorbox}[colback=navyblue!5, colframe=navyblue, width=\textwidth]
\centering
{\LARGE\bfseries\color{navyblue} BÁO CÁO BÀI TẬP LỚN}\\[0.4cm]
{\Large\bfseries Môn: CS112 -- Phân tích và Thiết kế Thuật toán}\\[0.6cm]
{\huge\bfseries\color{navyblue} Thuật toán Fourier-Motzkin\\[0.2cm]
Giải Quy hoạch Tuyến tính}
\end{tcolorbox}

\vspace{1.5cm}

\begin{tabular}{ll}
\textbf{Giảng viên hướng dẫn:} & \ldots\ldots\ldots\ldots\ldots\ldots\ldots \\[0.3cm]
\textbf{Sinh viên thực hiện:}   & \ldots\ldots\ldots\ldots\ldots\ldots\ldots \\[0.2cm]
\textbf{MSSV:}                  & \ldots\ldots\ldots\ldots\ldots\ldots\ldots \\[0.2cm]
\textbf{Lớp:}                   & \ldots\ldots\ldots\ldots\ldots\ldots\ldots \\
\end{tabular}

\vfill
{\large TP. Hồ Chí Minh, \the\year}
\end{titlepage}

% ── Mục lục ───────────────────────────────────────────────────────────────
\tableofcontents
\newpage

% ══════════════════════════════════════════════════════════════════════════
\section{Giới thiệu}
% ══════════════════════════════════════════════════════════════════════════

\subsection{Đặt vấn đề}

Quy hoạch tuyến tính (Linear Programming -- LP) là một trong những bài toán tối ưu hóa
quan trọng nhất trong khoa học máy tính và toán học ứng dụng. Bài toán LP xuất hiện
rộng rãi trong các lĩnh vực như kinh tế, logistics, lập lịch, phân bổ tài nguyên và
nhiều ứng dụng thực tiễn khác.

Bài tập lớn này tập trung vào việc cài đặt và trực quan hóa \textbf{thuật toán
Fourier-Motzkin} -- một trong những phương pháp cổ điển và nền tảng nhất để giải
bài toán LP thông qua phép chiếu (projection). Thuật toán được Joseph Fourier đề xuất
năm 1827 và Theodore Motzkin phát triển độc lập năm 1936.

\subsection{Mục tiêu}

\begin{itemize}[leftmargin=1.5cm]
    \item Nghiên cứu lý thuyết thuật toán Fourier-Motzkin và các tính chất toán học liên quan.
    \item Cài đặt thuật toán theo \textbf{hai hướng tiếp cận}: số học (algebraic) và hình học (geometric).
    \item Xây dựng ứng dụng web demo với giao diện trực quan, hỗ trợ giải trình từng bước.
    \item Phân tích độ phức tạp và đánh giá hiệu năng của thuật toán.
\end{itemize}

\subsection{Phạm vi}

Ứng dụng hỗ trợ giải bài toán LP dạng tổng quát với $n$ biến và $m$ ràng buộc
(dạng $\leq$, $\geq$, $=$). Phương pháp hình học chỉ áp dụng cho bài toán 2 biến
để có thể trực quan hóa trên mặt phẳng.

% ══════════════════════════════════════════════════════════════════════════
\section{Cơ sở lý thuyết}
% ══════════════════════════════════════════════════════════════════════════

\subsection{Bài toán Quy hoạch Tuyến tính}

\begin{definition}[Bài toán LP chuẩn tắc]
Bài toán quy hoạch tuyến tính có dạng:
\begin{equation}
    \begin{aligned}
        &\text{tối ưu hóa} \quad z = \vc^T \vx \\
        &\text{thỏa mãn} \quad \mA\vx \leq \vb \\
        &\phantom{\text{thỏa mãn}} \quad \vx \geq \mathbf{0}
    \end{aligned}
    \label{eq:lp_standard}
\end{equation}
trong đó $\vx \in \R^n$ là vector biến quyết định, $\vc \in \R^n$ là vector hệ số
hàm mục tiêu, $\mA \in \R^{m \times n}$ là ma trận ràng buộc, $\vb \in \R^m$ là
vector vế phải.
\end{definition}

\begin{definition}[Miền khả thi]
Tập hợp $\mathcal{F} = \{\vx \in \R^n \mid \mA\vx \leq \vb,\ \vx \geq \mathbf{0}\}$
được gọi là \textbf{miền khả thi} (feasible region). Mỗi điểm $\vx \in \mathcal{F}$
là một \textbf{nghiệm khả thi}.
\end{definition}

\begin{theorem}[Định lý điểm cực biên]
Nếu bài toán LP có nghiệm tối ưu hữu hạn, thì tồn tại ít nhất một \textbf{đỉnh}
(vertex) của miền khả thi $\mathcal{F}$ là nghiệm tối ưu.
\end{theorem}

\subsection{Thuật toán Fourier-Motzkin}

\subsubsection{Ý tưởng cốt lõi}

Thuật toán Fourier-Motzkin (FM) giải bài toán LP bằng cách \textbf{loại dần từng biến}
thông qua phép chiếu. Với mỗi biến $x_k$, thuật toán:

\begin{enumerate}
    \item Phân loại các ràng buộc thành ba nhóm theo hệ số của $x_k$:
    \begin{itemize}
        \item Nhóm $U$: $x_k \leq f_i(\vx_{-k})$ (hệ số dương)
        \item Nhóm $L$: $x_k \geq g_j(\vx_{-k})$ (hệ số âm)
        \item Nhóm $N$: không chứa $x_k$ (hệ số bằng 0)
    \end{itemize}
    \item Tạo ràng buộc mới bằng cách ghép mọi cặp $(l \in L,\ u \in U)$:
    \begin{equation}
        g_j(\vx_{-k}) \leq f_i(\vx_{-k})
        \label{eq:fm_pair}
    \end{equation}
    \item Hệ mới gồm nhóm $N$ và tất cả ràng buộc mới từ bước 2.
\end{enumerate}

\subsubsection{Thêm biến mục tiêu $z$}

Để tối ưu hóa hàm mục tiêu $z = \vc^T\vx$, ta thêm biến $z$ vào hệ:

\begin{itemize}
    \item \textbf{Bài toán max}: thêm ràng buộc $z \leq \vc^T\vx$, tức là
          $z - \vc^T\vx \leq 0$.
    \item \textbf{Bài toán min}: chuyển về $\max(-\vc^T\vx)$, thêm ràng buộc
          $z \leq -\vc^T\vx$, tức là $z + \vc^T\vx \leq 0$.
\end{itemize}

Sau khi loại hết $x_1, \ldots, x_n$, hệ còn lại chỉ chứa $z$, từ đó tìm được
$z^* = \min\{b_i/a_i \mid a_i > 0\}$.

\subsubsection{Thuật toán FM đầy đủ}

\begin{algorithm}[H]
\caption{Fourier-Motzkin LP Solver}
\label{alg:fm}
\begin{algorithmic}[1]
\Require $n$ biến, $m$ ràng buộc $\{a_i^T\vx \leq b_i\}$, hàm mục tiêu $\vc^T\vx$, loại (max/min)
\Ensure Nghiệm tối ưu $\vx^*$ và giá trị $z^*$, hoặc thông báo vô nghiệm/không bị chặn

\State \textbf{Bước 1: Chuẩn hóa} -- Đổi tất cả ràng buộc về dạng $\leq$
\For{mỗi ràng buộc $\geq$} nhân hai vế với $-1$ \EndFor
\For{mỗi ràng buộc $=$} tách thành hai ràng buộc $\leq$ \EndFor

\State \textbf{Bước 2: Thêm biến $z$}
\If{max} thêm $z - \vc^T\vx \leq 0$ \Else\ thêm $z + \vc^T\vx \leq 0$ \EndIf

\State \textbf{Bước 3: Loại biến FM}
\For{$k = 1$ \textbf{to} $n$}
    \State Phân loại ràng buộc: nhóm $U$, $L$, $N$ theo hệ số $x_k$
    \State Chuẩn hóa: $x_k \leq f_i$ (từ $U$), $x_k \geq g_j$ (từ $L$)
    \For{mỗi cặp $(g_j, f_i) \in L \times U$}
        \State Thêm ràng buộc mới: $(g_j\text{-coeffs} - f_i\text{-coeffs}) \cdot \vx \leq f_i\text{-rhs} - g_j\text{-rhs}$
    \EndFor
    \State Hệ mới $\leftarrow$ nhóm $N$ $\cup$ các ràng buộc mới
\EndFor

\State \textbf{Bước 4: Tìm $z^*$}
\State Thu thập tất cả ràng buộc dạng $a_z \cdot z \leq b$ (chỉ chứa $z$)
\State $z^* \leftarrow \min\{b/a_z \mid a_z > 0\}$
\If{không có chặn trên} \Return \textbf{Không bị chặn} \EndIf

\State \textbf{Bước 5: Back substitution}
\For{$k = n$ \textbf{downto} $1$}
    \State Thế $z = z^*$ và $x_{k+1}, \ldots, x_n$ đã biết vào hệ tại bước $k$
    \State Tìm khoảng $[l_k, u_k]$ cho $x_k$, chọn $x_k \in [l_k, u_k]$
\EndFor
\State \Return $\vx^* = (x_1^*, \ldots, x_n^*)$, $z^*$
\end{algorithmic}
\end{algorithm}

\subsubsection{Độ phức tạp}

\begin{theorem}[Độ phức tạp Fourier-Motzkin]
Sau khi loại biến $x_k$, số ràng buộc có thể tăng từ $m$ lên tối đa
$\lfloor m^2/4 \rfloor + m_N$ (do ghép $|L| \times |U|$ cặp, tối đa khi $|L| = |U| = m/2$).
Sau $n$ bước loại biến, số ràng buộc có thể đạt $O(m^{2^n})$ trong trường hợp xấu nhất.
\end{theorem}

\begin{remark}
Mặc dù độ phức tạp lý thuyết là hàm mũ kép, trong thực tế nhiều ràng buộc bị loại
(redundant) hoặc hệ số bằng 0, nên thuật toán vẫn hiệu quả với các bài toán vừa và nhỏ.
Đây là lý do FM được dùng rộng rãi trong kiểm tra tính khả thi và phân tích hệ bất đẳng thức.
\end{remark}

\subsection{Phương pháp hình học (2 biến)}

Với bài toán 2 biến, miền khả thi $\mathcal{F}$ là một \textbf{đa giác lồi} (convex polygon)
trên mặt phẳng $Ox_1x_2$. Phương pháp hình học gồm 4 bước:

\begin{enumerate}
    \item \textbf{Vẽ đường biên}: Mỗi ràng buộc $a_1x_1 + a_2x_2 \leq b$ xác định
          một nửa mặt phẳng. Đường biên là đường thẳng $a_1x_1 + a_2x_2 = b$.
    \item \textbf{Xác định miền khả thi}: Giao của tất cả các nửa mặt phẳng.
          Tìm các đỉnh bằng cách giải hệ $2 \times 2$ từ mọi cặp ràng buộc.
    \item \textbf{Tính $z$ tại các đỉnh}: Theo Định lý điểm cực biên, nghiệm tối ưu
          nằm tại một đỉnh của $\mathcal{F}$.
    \item \textbf{Đường mức}: Tịnh tiến đường mức $z = \vc^T\vx = c$ theo hướng
          gradient $\nabla z = \vc$ (max) hoặc $-\vc$ (min) cho đến khi chạm biên.
\end{enumerate}

\begin{figure}[H]
\centering
\begin{tikzpicture}
\begin{axis}[
    width=10cm, height=8cm,
    xlabel={$x_1$}, ylabel={$x_2$},
    xmin=-0.3, xmax=4.5, ymin=-0.3, ymax=5,
    grid=both, grid style={line width=0.3pt, draw=gray!30},
    axis lines=left,
    tick label style={font=\small},
    legend pos=north east,
    legend style={font=\small}
]
% Miền khả thi (tô màu)
\addplot[fill=navyblue!15, draw=none, opacity=0.5]
    coordinates {(0,0) (3,0) (2,2) (0,4)} -- cycle;

% Đường biên
\addplot[thick, navyblue, domain=0:4, samples=2] {4-x}
    node[right, font=\small] {$x_1+x_2=4$};
\addplot[thick, red!70!black, domain=0:3, samples=2] {6-2*x}
    node[right, font=\small] {$2x_1+x_2=6$};
\addplot[thick, gray, domain=0:0.1, samples=2] coordinates {(0,0)(0,4.5)}
    node[above, font=\small] {$x_1=0$};
\addplot[thick, gray, domain=0:4.5, samples=2] {0}
    node[right, font=\small] {$x_2=0$};

% Đường mức tối ưu
\addplot[thick, dashed, codegreen, domain=0:4, samples=2] {(10-3*x)/2}
    node[right, font=\small, color=codegreen] {$z^*=10$};

% Các đỉnh
\addplot[only marks, mark=*, mark size=3pt, navyblue]
    coordinates {(0,0) (3,0) (2,2) (0,4)};
\node[above left, font=\small] at (axis cs:0,0) {$O(0,0)$};
\node[below right, font=\small] at (axis cs:3,0) {$B(3,0)$};
\node[above right, font=\small, color=red!70!black] at (axis cs:2,2) {$\mathbf{C^*(2,2)}$};
\node[above right, font=\small] at (axis cs:0,4) {$D(0,4)$};

% Mũi tên gradient
\draw[->, thick, codegreen!80!black, line width=1.5pt]
    (axis cs:1,1) -- (axis cs:1.6,1.4)
    node[right, font=\small, color=codegreen!80!black] {$\nabla z$};

\end{axis}
\end{tikzpicture}
\caption{Minh họa phương pháp hình học: bài toán $\max z = 3x_1 + 2x_2$
với $x_1+x_2 \leq 4$, $2x_1+x_2 \leq 6$, $x_1,x_2 \geq 0$.
Nghiệm tối ưu tại $C^*(2,2)$ với $z^* = 10$.}
\label{fig:geometric}
\end{figure}

% ══════════════════════════════════════════════════════════════════════════
\section{Thiết kế hệ thống}
% ══════════════════════════════════════════════════════════════════════════

\subsection{Kiến trúc tổng thể}

Ứng dụng được xây dựng theo mô hình \textbf{Client-Server} với Flask (Python) làm backend
và HTML/CSS/JavaScript thuần làm frontend.

\begin{figure}[H]
\centering
\begin{tikzpicture}[
    node distance=1.5cm,
    box/.style={rectangle, draw=navyblue, fill=lightblue, rounded corners=4pt,
                minimum width=3cm, minimum height=1cm, font=\small\bfseries,
                text centered},
    arrow/.style={->, thick, navyblue}
]
\node[box, fill=green!10, draw=codegreen] (browser) {Trình duyệt\\(HTML/CSS/JS)};
\node[box, right=3cm of browser] (flask) {Flask App\\(main.py)};
\node[box, below=1.2cm of flask] (solver) {Solver Module};
\node[box, below left=0.8cm and -0.5cm of solver] (core) {core\_exact.py\\(FM số học)};
\node[box, below right=0.8cm and -0.5cm of solver] (geo) {geometric.py\\(FM hình học)};
\node[box, below=1.2cm of core] (exact) {exact.py\\(SymPy)};
\node[box, below=1.2cm of geo] (reason) {reasoning.py\\(Ghi bước)};
\node[box, right=1.5cm of solver] (parser) {parser.py\\(Parse text)};

\draw[arrow] (browser) -- node[above, font=\tiny] {POST /api/solve} (flask);
\draw[arrow] (flask) -- node[right, font=\tiny] {} (solver);
\draw[arrow] (solver) -- (core);
\draw[arrow] (solver) -- (geo);
\draw[arrow] (core) -- (exact);
\draw[arrow] (core) -- (reason);
\draw[arrow] (geo) -- (reason);
\draw[arrow] (flask) -- (parser);
\draw[arrow, dashed] (flask) to[bend left=20] node[above, font=\tiny] {JSON response} (browser);
\end{tikzpicture}
\caption{Kiến trúc hệ thống Fourier-Motzkin LP Solver}
\label{fig:architecture}
\end{figure}

\subsection{Cấu trúc thư mục}

\begin{lstlisting}[language=bash, caption={Cấu trúc dự án}, label={lst:structure}]
fourier-motzkin/
|-- main.py              # Flask app, API endpoints
|-- requirements.txt     # Dependencies (Flask, SymPy, ...)
|-- Procfile             # Gunicorn deploy config
|-- solver/
|   |-- __init__.py
|   |-- core.py          # FM solver (float, legacy)
|   |-- core_exact.py    # FM solver (SymPy exact, chinh)
|   |-- exact.py         # Tien ich so hoc exact (SymPy)
|   |-- geometric.py     # Phuong phap hinh hoc (2 bien)
|   |-- parser.py        # Parse rang buoc dang text
|   `-- reasoning.py     # Ghi lai cac buoc lap luan
|-- templates/
|   `-- index.html       # Giao dien web (SPA)
`-- static/
    `-- style.css        # CSS bo sung
\end{lstlisting}

\subsection{Mô tả các module}

\subsubsection{Module \texttt{exact.py} -- Số học chính xác}

Module này sử dụng thư viện \textbf{SymPy} để thực hiện tính toán với số hữu tỉ chính xác,
tránh sai số dấu phẩy động. Các hàm chính:

\begin{itemize}
    \item \texttt{to\_exact\_number(value)}: Chuyển đổi giá trị bất kỳ (int, float, str)
          thành số SymPy exact (\texttt{Rational}, \texttt{Integer}).
    \item \texttt{fmt\_exact(value)}: Định dạng số exact thành chuỗi hiển thị.
    \item \texttt{constraint\_to\_exact(c)}: Chuyển ràng buộc sang dạng exact.
    \item \texttt{format\_linear\_expression(coeffs, var\_names)}: Định dạng biểu thức tuyến tính.
\end{itemize}

\begin{lstlisting}[caption={Ví dụ sử dụng exact.py}, label={lst:exact}]
from solver.exact import to_exact_number, fmt_exact
from sympy import Rational

# Chuyen float thanh Rational chinh xac
x = to_exact_number(0.333333)
# => Rational('1/3') thay vi 0.333333...

# Tinh toan chinh xac
a = to_exact_number("1/3")
b = to_exact_number("2/5")
print(fmt_exact(a + b))  # => "11/15"
\end{lstlisting}

\subsubsection{Module \texttt{core\_exact.py} -- FM số học}

Đây là module chính cài đặt thuật toán Fourier-Motzkin với số học exact. Pipeline:

\begin{enumerate}
    \item \texttt{chuan\_hoa\_he()}: Chuẩn hóa tất cả ràng buộc về dạng $\leq$.
    \item \texttt{them\_bien\_z()}: Thêm biến $z$ vào hệ.
    \item \texttt{co\_lap\_bien()}: Phân loại ràng buộc theo hệ số của biến cần loại.
    \item \texttt{ghep\_cap\_FM()}: Ghép cặp $(L, U)$ tạo ràng buộc mới.
    \item \texttt{tim\_z\_max()}: Tìm chặn trên của $z$.
    \item \texttt{back\_substitute\_FM()}: Thế ngược tìm nghiệm.
    \item \texttt{kiem\_tra\_nghiem()}: Kiểm tra nghiệm thỏa mãn hệ gốc.
\end{enumerate}

\subsubsection{Module \texttt{geometric.py} -- FM hình học}

Cài đặt phương pháp hình học cho bài toán 2 biến:

\begin{itemize}
    \item \texttt{phan\_loai\_duong\_bien()}: Phân loại đường biên (trục, đứng, ngang, thường).
    \item \texttt{giai\_he\_phuong\_trinh()}: Giải hệ $2 \times 2$ tìm giao điểm.
    \item \texttt{tinh\_huong\_vec()}: Tính vector hướng vào miền nghiệm.
    \item \texttt{sap\_xep\_dinh\_convex()}: Sắp xếp đỉnh theo thứ tự convex hull.
    \item \texttt{giai\_hinh\_hoc()}: Hàm chính giải bài toán hình học.
\end{itemize}

\subsubsection{Module \texttt{parser.py} -- Parse ràng buộc text}

Cho phép người dùng nhập ràng buộc dạng văn bản tự nhiên, hỗ trợ:
\begin{itemize}
    \item Hệ số là biểu thức toán học: \texttt{2\^{}3}, \texttt{1/3}, \texttt{sqrt(2)}.
    \item Hàm lượng giác trong hệ số: \texttt{sin(pi/6)}, \texttt{cos(pi/3)}.
    \item Phát hiện và báo lỗi chi tiết với số dòng và gợi ý sửa.
\end{itemize}

\subsubsection{Module \texttt{reasoning.py} -- Ghi bước lập luận}

\texttt{ReasoningEngine} ghi lại từng bước tính toán với 4 loại:
\texttt{INFO}, \texttt{TINH\_TOAN}, \texttt{KET\_LUAN}, \texttt{CANH\_BAO}.
Kết quả được trả về dạng JSON để hiển thị trên giao diện.

% ══════════════════════════════════════════════════════════════════════════
\section{Cài đặt chi tiết}
% ══════════════════════════════════════════════════════════════════════════

\subsection{Hàm chuẩn hóa ràng buộc}

\begin{lstlisting}[caption={Chuẩn hóa ràng buộc về dạng $\leq$}, label={lst:normalize}]
def chuan_hoa_he(constraints, obj_coeffs, obj_type, engine):
    """Chuẩn hoá về dạng ≤."""
    result = []
    for c in constraints:
        c = constraint_to_exact(c)
        if c["sense"] == "<=":
            result.append({"coeffs": c["coeffs"][:], "rhs": c["rhs"], "sense": "<="})
        elif c["sense"] == ">=":
            # Nhân hai vế với -1: a*x >= b  =>  -a*x <= -b
            nc = {"coeffs": [-v for v in c["coeffs"]], "rhs": -c["rhs"], "sense": "<="}
            result.append(nc)
        elif c["sense"] == "=":
            # Tách thành 2 bất đẳng thức: a*x = b  =>  a*x <= b  và  -a*x <= -b
            cle = {"coeffs": c["coeffs"][:], "rhs": c["rhs"], "sense": "<="}
            cge = {"coeffs": [-v for v in c["coeffs"]], "rhs": -c["rhs"], "sense": "<="}
            result.extend([cle, cge])
    return result
\end{lstlisting}

\subsection{Hàm phân loại và ghép cặp FM}

\begin{lstlisting}[caption={Phân loại ràng buộc và ghép cặp FM}, label={lst:fm_core}]
def co_lap_bien(constraints, var_idx):
    """Phân loại ràng buộc theo hệ số của biến var_idx."""
    nhom_U, nhom_L, nhom_N = [], [], []
    for c in constraints:
        a = c["coeffs"][var_idx]
        if is_zero(a):
            nhom_N.append(c)
            continue
        # Chuẩn hóa: chia hai vế cho |a|
        sc = [-(c["coeffs"][j]) / a for j in range(len(c["coeffs"]))]
        sc[var_idx] = to_exact_number(0)
        sr = c["rhs"] / a
        bc = {"coeffs": sc, "rhs": sr, "sense": "<=" if is_positive(a) else ">="}
        (nhom_U if is_positive(a) else nhom_L).append(bc)
    return nhom_U, nhom_L, nhom_N

def ghep_cap_FM(nhom_L, nhom_U, nhom_N, var_idx, var_names, engine):
    """Ghép mọi cặp (l in L, u in U): f_L <= f_U."""
    result = list(nhom_N)
    for l in nhom_L:
        for u in nhom_U:
            # Ràng buộc mới: (l_coeffs - u_coeffs)*x <= u_rhs - l_rhs
            new_coeffs = [lc - uc for lc, uc in zip(l["coeffs"], u["coeffs"])]
            new_rhs = u["rhs"] - l["rhs"]
            result.append({"coeffs": new_coeffs, "rhs": new_rhs, "sense": "<="})
    return result
\end{lstlisting}

\subsection{Back substitution}

Sau khi tìm được $z^*$, thuật toán thế ngược từ $x_n$ về $x_1$:

\begin{lstlisting}[caption={Back substitution tìm nghiệm}, label={lst:backsubst}]
def back_substitute_FM(steps_constraints, z_max, n, engine):
    """Thế z=z_max và các biến đã biết, tìm từng xi từ xn về x1."""
    solution = [to_exact_number(0)] * n
    for k in range(n - 1, -1, -1):
        cons = steps_constraints[k]  # Hệ trước khi loại x_k
        lower, upper = None, None
        for c in cons:
            a = c["coeffs"][k]
            if is_zero(a): continue
            # Thế z và x_{k+1},...,x_{n-1} đã biết
            rhs_val = c["rhs"]
            rhs_val -= c["coeffs"][n] * z_max  # trừ phần z
            for j in range(k + 1, n):
                rhs_val -= c["coeffs"][j] * solution[j]
            bound = rhs_val / a
            if is_positive(a):
                upper = bound if upper is None or is_less(bound, upper) else upper
            else:
                lower = bound if lower is None or is_greater(bound, lower) else lower
        # Chọn giá trị trong [lower, upper]
        if lower is not None and upper is not None:
            val = lower  # hoặc (lower+upper)/2
        elif lower is not None:
            val = lower
        elif upper is not None:
            val = max(upper, to_exact_number(0))
        else:
            val = to_exact_number(0)
        solution[k] = val
    return solution
\end{lstlisting}

\subsection{Phương pháp hình học -- Tìm đỉnh khả thi}

\begin{lstlisting}[caption={Tìm các đỉnh của miền khả thi}, label={lst:vertices}]
def giai_hinh_hoc(n, constraints, obj_coeffs, obj_type):
    """Giải LP 2 biến bằng phương pháp hình học."""
    from itertools import combinations
    feasible_vertices = []
    seen = set()
    # Xét mọi cặp ràng buộc, giải hệ 2x2
    for (i, j) in combinations(range(len(constraints)), 2):
        sol = giai_he_phuong_trinh(constraints[i], constraints[j])
        if sol is None: continue
        key = (round(sol[0], 6), round(sol[1], 6))
        if key in seen: continue
        seen.add(key)
        # Kiểm tra điểm có thỏa mãn TẤT CẢ ràng buộc không
        if not _kiem_tra(sol, constraints): continue
        feasible_vertices.append({"point": list(sol), ...})
    # Tính z tại mỗi đỉnh, chọn đỉnh tối ưu
    best_z = None
    for v in feasible_vertices:
        z = sum(obj_coeffs[k] * v["point"][k] for k in range(2))
        if best_z is None or (obj_type=="max" and z > best_z):
            best_z = z
    return {"feasible": True, "z": best_z, ...}
\end{lstlisting}

% ══════════════════════════════════════════════════════════════════════════
\section{Giao diện người dùng (UI Demo)}
% ══════════════════════════════════════════════════════════════════════════

\subsection{Tổng quan giao diện}

Giao diện web được xây dựng bằng HTML/CSS/JavaScript thuần (không dùng framework),
với thiết kế hiện đại, responsive. Màu chủ đạo là navy blue (\texttt{\#1a237e}).

\subsection{Các thành phần chính}

\subsubsection{Thanh công cụ (Toolbar)}

\begin{itemize}
    \item \textbf{Số biến}: Nhập số biến $n \geq 2$, nhấn ``Áp dụng'' để cập nhật form.
    \item \textbf{Mục tiêu}: Toggle MAX/MIN.
    \item \textbf{Ví dụ mẫu / Reset}: Nạp bài toán mẫu hoặc xóa toàn bộ.
\end{itemize}

\subsubsection{Nhập hàm mục tiêu và ràng buộc}

Hỗ trợ \textbf{hai chế độ nhập}:
\begin{itemize}
    \item \textbf{Form}: Nhập từng hệ số vào ô input, chọn dấu ($\leq$, $\geq$, $=$),
          nhập vế phải. Có nút thêm/xóa ràng buộc động.
    \item \textbf{Text}: Nhập tự do dạng văn bản, ví dụ:
          \texttt{2x1 + x2 <= 6}, \texttt{sin(pi/6)x1 + sqrt(2)x2 >= 1}.
          Parser tự động phân tích và báo lỗi chi tiết.
\end{itemize}

Checkbox ``Tự động thêm điều kiện không âm $x_i \geq 0$'' giúp tiết kiệm thao tác.

\subsubsection{Chọn phương pháp giải}

Toggle giữa hai phương pháp:
\begin{itemize}
    \item \textbf{Số học (FM)}: Áp dụng cho mọi $n \geq 1$.
    \item \textbf{Hình học}: Chỉ áp dụng cho $n = 2$, hiển thị biểu đồ 2D.
\end{itemize}

\subsubsection{Kết quả và các tab}

Sau khi nhấn ``GIẢI'', kết quả hiển thị qua 4 tab:

\begin{enumerate}
    \item \textbf{Các bước FM}: Accordion hiển thị hệ ràng buộc sau mỗi bước loại biến.
          Ràng buộc mới được tô xanh (★), ràng buộc bị loại gạch ngang.
    \item \textbf{Biểu đồ}: Chart.js vẽ miền khả thi, các đỉnh (có nhãn tọa độ),
          đường biên, đường mức tối ưu và mũi tên gradient. Chỉ hiển thị khi $n = 2$.
    \item \textbf{Reasoning}: Hiển thị từng bước lập luận chi tiết với màu sắc
          phân loại (INFO, TÍNH TOÁN, KẾT LUẬN, CẢNH BÁO). Có bộ lọc theo loại.
    \item \textbf{Xuất file}: Tải về file \texttt{.txt} chứa toàn bộ giải trình
          (cả hai phương pháp đều hỗ trợ).
\end{enumerate}

\subsection{Luồng xử lý API}

\begin{figure}[H]
\centering
\begin{tikzpicture}[
    node distance=1.2cm,
    step/.style={rectangle, draw=navyblue, fill=lightblue, rounded corners=3pt,
                 minimum width=4.5cm, minimum height=0.8cm, font=\small, text centered},
    decision/.style={diamond, draw=navyblue, fill=yellow!20, font=\small,
                     minimum width=2.5cm, minimum height=0.8cm, text centered, aspect=2},
    arrow/.style={->, thick, navyblue}
]
\node[step] (input) {Người dùng nhập bài toán};
\node[step, below of=input] (collect) {collectInput() -- Thu thập dữ liệu};
\node[step, below of=collect] (post) {POST /api/solve (JSON)};
\node[decision, below of=post, yshift=-0.3cm] (method) {Phương pháp?};
\node[step, below left=1cm and 1.5cm of method] (fm) {giai\_bai\_toan()\\(FM số học)};
\node[step, below right=1cm and 1.5cm of method] (geo) {giai\_hinh\_hoc()\\(Hình học)};
\node[step, below=2.5cm of method] (json) {Trả về JSON kết quả};
\node[step, below of=json] (render) {Render kết quả + biểu đồ};

\draw[arrow] (input) -- (collect);
\draw[arrow] (collect) -- (post);
\draw[arrow] (post) -- (method);
\draw[arrow] (method) -- node[above left, font=\tiny] {algebraic} (fm);
\draw[arrow] (method) -- node[above right, font=\tiny] {geometric} (geo);
\draw[arrow] (fm) -- (json);
\draw[arrow] (geo) -- (json);
\draw[arrow] (json) -- (render);
\end{tikzpicture}
\caption{Luồng xử lý từ input đến kết quả}
\label{fig:flow}
\end{figure}

\subsection{Công nghệ sử dụng}

\begin{table}[H]
\centering
\caption{Công nghệ và thư viện sử dụng}
\label{tab:tech}
\begin{tabularx}{\textwidth}{lXl}
\toprule
\textbf{Thành phần} & \textbf{Mô tả} & \textbf{Phiên bản} \\
\midrule
Python & Ngôn ngữ lập trình chính & 3.x \\
Flask & Web framework (backend) & 3.1.3 \\
SymPy & Số học chính xác (Rational, Integer) & 1.14.0 \\
Chart.js & Vẽ biểu đồ 2D (CDN) & 4.4.1 \\
Gunicorn & WSGI server (deploy) & 25.3.0 \\
HTML/CSS/JS & Frontend (không dùng framework) & -- \\
\bottomrule
\end{tabularx}
\end{table}

% ══════════════════════════════════════════════════════════════════════════
\section{Ví dụ minh họa}
% ══════════════════════════════════════════════════════════════════════════

\subsection{Ví dụ 1: Bài toán 2 biến (MAX)}

\begin{infobox}[Bài toán]
\begin{align*}
    \max\quad & z = 3x_1 + 2x_2 \\
    \text{s.t.}\quad & x_1 + x_2 \leq 4 \\
    & 2x_1 + x_2 \leq 6 \\
    & x_1, x_2 \geq 0
\end{align*}
\end{infobox}

\subsubsection{Giải bằng phương pháp số học (FM)}

\textbf{Bước 1: Chuẩn hóa và thêm biến $z$}

Hệ ràng buộc sau chuẩn hóa (đã có dạng $\leq$):
\begin{align}
    x_1 + x_2 &\leq 4 \label{c1}\\
    2x_1 + x_2 &\leq 6 \label{c2}\\
    -x_1 &\leq 0 \label{c3}\\
    -x_2 &\leq 0 \label{c4}
\end{align}

Thêm biến $z$: $z - 3x_1 - 2x_2 \leq 0$ \hfill (ràng buộc $z \leq 3x_1 + 2x_2$)

\textbf{Bước 2: Loại $x_1$}

Phân loại theo hệ số $x_1$:
\begin{itemize}
    \item Nhóm $U$ ($x_1 \leq \ldots$): từ \eqref{c1}: $x_1 \leq 4 - x_2$;
          từ \eqref{c2}: $x_1 \leq 3 - \frac{1}{2}x_2$;
          từ ràng buộc $z$: $x_1 \leq \frac{1}{3}(z - 2x_2) \cdot (-1)$... (sau chuẩn hóa)
    \item Nhóm $L$ ($x_1 \geq \ldots$): từ \eqref{c3}: $x_1 \geq 0$
    \item Nhóm $N$: \eqref{c4}
\end{itemize}

Ghép cặp $L \times U$ tạo ràng buộc mới (loại $x_1$):
\begin{align}
    0 &\leq 4 - x_2 \implies x_2 \leq 4 \\
    0 &\leq 3 - \tfrac{1}{2}x_2 \implies x_2 \leq 6 \\
    0 &\leq \tfrac{1}{3}(z - 2x_2) \cdot (-1) \implies \ldots
\end{align}

\textbf{Bước 3: Loại $x_2$} (tương tự)

\textbf{Bước 4: Tìm $z^*$}

Sau khi loại hết $x_1, x_2$, thu được: $z \leq 10$

$$\boxed{z^* = 10}$$

\textbf{Bước 5: Back substitution}

Thế $z = 10$ vào hệ tại bước loại $x_2$: tìm được $x_2 = 2$.\\
Thế $z = 10$, $x_2 = 2$ vào hệ tại bước loại $x_1$: tìm được $x_1 = 2$.

\begin{resultbox}[Kết quả]
$$x_1^* = 2,\quad x_2^* = 2,\quad z^* = 3(2) + 2(2) = \mathbf{10}$$
\end{resultbox}

\subsubsection{Giải bằng phương pháp hình học}

Các đỉnh của miền khả thi (giao điểm các cặp ràng buộc):

\begin{table}[H]
\centering
\caption{Giá trị $z$ tại các đỉnh -- Ví dụ 1}
\label{tab:ex1_vertices}
\begin{tabular}{cccc}
\toprule
\textbf{Đỉnh} & $x_1$ & $x_2$ & $z = 3x_1 + 2x_2$ \\
\midrule
$O$ & 0 & 0 & 0 \\
$B$ & 3 & 0 & 9 \\
$C^*$ & 2 & 2 & \textbf{10} $\leftarrow$ MAX \\
$D$ & 0 & 4 & 8 \\
\bottomrule
\end{tabular}
\end{table}

Nghiệm tối ưu tại $C^*(2, 2)$ với $z^* = 10$. ✓

\subsection{Ví dụ 2: Bài toán 3 biến (MAX)}

\begin{infobox}[Bài toán]
\begin{align*}
    \max\quad & z = 2x_1 + 3x_2 + x_3 \\
    \text{s.t.}\quad & x_1 + x_2 + x_3 \leq 6 \\
    & x_2 + x_3 \leq 4 \\
    & x_1 + x_2 \leq 4 \\
    & x_1, x_2, x_3 \geq 0
\end{align*}
\end{infobox}

Bài toán 3 biến không thể giải bằng phương pháp hình học (cần không gian 3D).
Áp dụng FM số học:

\begin{resultbox}[Kết quả]
$$x_1^* = 2,\quad x_2^* = 2,\quad x_3^* = 2,\quad z^* = 2(2)+3(2)+1(2) = \mathbf{12}$$
\end{resultbox}

\subsection{Ví dụ 3: Bài toán MIN với hệ số phân số}

\begin{infobox}[Bài toán]
\begin{align*}
    \min\quad & z = -x_1 - 2x_2 \\
    \text{s.t.}\quad & 2x_1 + x_2 \leq 6 \\
    & x_1 + x_2 \geq 2 \\
    & -x_1 + x_2 \geq 3 \\
    & x_1, x_2 \geq 0
\end{align*}
\end{infobox}

Bài toán MIN được chuyển về $\max(-z) = \max(x_1 + 2x_2)$.

\begin{resultbox}[Kết quả]
Nghiệm tối ưu: $x_1^* = 0$, $x_2^* = 6$, $z^* = -12$
\end{resultbox}

% ══════════════════════════════════════════════════════════════════════════
\section{Phân tích và đánh giá}
% ══════════════════════════════════════════════════════════════════════════

\subsection{Phân tích độ phức tạp}

\subsubsection{Độ phức tạp thời gian}

\begin{table}[H]
\centering
\caption{Phân tích độ phức tạp theo số biến và ràng buộc}
\label{tab:complexity}
\begin{tabular}{lll}
\toprule
\textbf{Bước} & \textbf{Độ phức tạp} & \textbf{Ghi chú} \\
\midrule
Chuẩn hóa & $O(m)$ & $m$ ràng buộc \\
Thêm biến $z$ & $O(m)$ & Thêm 1 ràng buộc \\
Loại 1 biến & $O(m^2)$ & Ghép $|L| \times |U| \leq m^2/4$ cặp \\
Loại $n$ biến & $O(m^{2^n})$ & Trường hợp xấu nhất \\
Tìm $z^*$ & $O(m')$ & $m'$ = số ràng buộc còn lại \\
Back substitution & $O(n \cdot m)$ & $n$ bước, mỗi bước $O(m)$ \\
\midrule
\textbf{Tổng} & $O(m^{2^n})$ & Hàm mũ kép theo $n$ \\
\bottomrule
\end{tabular}
\end{table}

\subsubsection{Độ phức tạp không gian}

Số ràng buộc sau $k$ bước loại biến tối đa là $O(m^{2^k})$, dẫn đến không gian
lưu trữ cũng là $O(m^{2^n})$ trong trường hợp xấu nhất.

\subsubsection{So sánh với Simplex}

\begin{table}[H]
\centering
\caption{So sánh Fourier-Motzkin và Simplex}
\label{tab:compare}
\begin{tabularx}{\textwidth}{lXX}
\toprule
\textbf{Tiêu chí} & \textbf{Fourier-Motzkin} & \textbf{Simplex} \\
\midrule
Độ phức tạp (xấu nhất) & $O(m^{2^n})$ -- hàm mũ kép & $O(2^n)$ -- hàm mũ \\
Độ phức tạp (thực tế) & Tốt với $n$ nhỏ & Đa thức trong hầu hết trường hợp \\
Tính chính xác & Exact (với SymPy) & Có thể có sai số số học \\
Giải trình bước & Rõ ràng, dễ hiểu & Phức tạp hơn (pivot) \\
Ứng dụng chính & Kiểm tra khả thi, $n$ nhỏ & Bài toán lớn thực tế \\
Phát hiện vô nghiệm & Tự nhiên trong quá trình & Cần phase I \\
\bottomrule
\end{tabularx}
\end{table}

\subsection{Ưu điểm của cài đặt}

\begin{enumerate}
    \item \textbf{Số học chính xác}: Sử dụng SymPy \texttt{Rational} thay vì float,
          tránh hoàn toàn sai số tích lũy. Ví dụ: $1/3 + 2/3 = 1$ (không phải $0.9999...$).

    \item \textbf{Hỗ trợ hệ số phức tạp}: Parser cho phép nhập hệ số dạng
          $\sin(\pi/6)$, $\sqrt{2}$, $2^3$, $1/3$ -- được tính chính xác bằng SymPy.

    \item \textbf{Giải trình chi tiết}: Mỗi bước tính toán được ghi lại với
          \texttt{ReasoningEngine}, giúp người học hiểu rõ quá trình.

    \item \textbf{Hai phương pháp bổ sung nhau}: FM số học cho kết quả chính xác
          với mọi $n$; hình học cho trực quan hóa với $n = 2$.

    \item \textbf{Giao diện thân thiện}: Responsive, hỗ trợ cả nhập form và text,
          có biểu đồ tương tác, xuất file giải trình.
\end{enumerate}

\subsection{Hạn chế}

\begin{enumerate}
    \item \textbf{Bùng nổ ràng buộc}: Với $n$ lớn (> 5-6 biến), số ràng buộc
          tăng theo hàm mũ kép, làm chậm đáng kể.

    \item \textbf{Hình học chỉ cho $n = 2$}: Không thể trực quan hóa không gian
          3 chiều trở lên trên giao diện 2D.

    \item \textbf{Back substitution}: Trong một số trường hợp đặc biệt (miền nghiệm
          không bị chặn theo một chiều), việc chọn giá trị trong khoảng $[l, u]$
          có thể không cho nghiệm tối ưu duy nhất.
\end{enumerate}

% ══════════════════════════════════════════════════════════════════════════
\section{Hướng dẫn sử dụng}
% ══════════════════════════════════════════════════════════════════════════

\subsection{Cài đặt và chạy}

\begin{lstlisting}[language=bash, caption={Cài đặt và chạy ứng dụng}]
# 1. Clone hoặc tải source code
cd fourier-motzkin

# 2. Tao virtual environment
python -m venv .venv
.venv\Scripts\activate   # Windows

# 3. Cai dependencies
pip install -r requirements.txt

# 4. Chay ung dung
python main.py
# => Truy cap http://localhost:5000
\end{lstlisting}

\subsection{Sử dụng API trực tiếp}

\begin{lstlisting}[language=bash, caption={Gọi API /api/solve}]
curl -X POST http://localhost:5000/api/solve \
  -H "Content-Type: application/json" \
  -d '{
    "n": 2,
    "obj_coeffs": [3, 2],
    "obj_type": "max",
    "method": "algebraic",
    "input_mode": "form",
    "constraints": [
      {"coeffs": [1, 1], "sense": "<=", "rhs": 4},
      {"coeffs": [2, 1], "sense": "<=", "rhs": 6},
      {"coeffs": [-1, 0], "sense": "<=", "rhs": 0},
      {"coeffs": [0, -1], "sense": "<=", "rhs": 0}
    ]
  }'
\end{lstlisting}

\textbf{Response JSON} trả về:
\begin{itemize}
    \item \texttt{feasible}: \texttt{true/false}
    \item \texttt{solution}: \texttt{\{"x1": "2", "x2": "2"\}}
    \item \texttt{z}: \texttt{"10"}
    \item \texttt{steps\_constraints}: Danh sách hệ ràng buộc sau mỗi bước FM
    \item \texttt{reasoning}: Chi tiết từng bước lập luận
    \item \texttt{vertices}, \texttt{boundary\_lines}, \texttt{level\_lines}: Dữ liệu biểu đồ
\end{itemize}

% ══════════════════════════════════════════════════════════════════════════
\section{Kết luận}
% ══════════════════════════════════════════════════════════════════════════

\subsection{Tổng kết}

Bài tập lớn đã hoàn thành các mục tiêu đề ra:

\begin{enumerate}
    \item \textbf{Nghiên cứu lý thuyết}: Nắm vững thuật toán Fourier-Motzkin,
          cơ sở toán học (phép chiếu, ghép cặp bất đẳng thức), và phương pháp hình học.

    \item \textbf{Cài đặt FM số học} (\texttt{core\_exact.py}): Pipeline đầy đủ
          gồm chuẩn hóa, thêm biến $z$, loại biến FM, tìm $z^*$, back substitution,
          kiểm tra nghiệm. Sử dụng SymPy để đảm bảo tính chính xác.

    \item \textbf{Cài đặt FM hình học} (\texttt{geometric.py}): Tìm đỉnh miền khả thi,
          tính $z$ tại các đỉnh, xác định nghiệm tối ưu, vẽ đường mức.

    \item \textbf{Xây dựng UI demo}: Ứng dụng web Flask với giao diện trực quan,
          hỗ trợ nhập liệu linh hoạt, hiển thị biểu đồ Chart.js, giải trình từng bước,
          xuất file kết quả.

    \item \textbf{Tính năng nổi bật}: Parser hỗ trợ hệ số phức tạp ($\sin$, $\sqrt{}$,
          phân số), số học exact với SymPy, reasoning engine ghi lại mọi bước.
\end{enumerate}

\subsection{Hướng phát triển}

\begin{itemize}
    \item Tối ưu hóa FM bằng cách loại bỏ ràng buộc dư thừa (redundant constraints)
          trước mỗi bước để giảm bùng nổ số ràng buộc.
    \item Thêm trực quan hóa 3D cho bài toán 3 biến (WebGL/Three.js).
    \item Cài đặt thêm thuật toán Simplex để so sánh trực tiếp.
    \item Hỗ trợ bài toán LP nguyên (Integer LP) bằng Branch and Bound.
    \item Thêm tính năng lưu/tải bài toán, chia sẻ link kết quả.
\end{itemize}

% ══════════════════════════════════════════════════════════════════════════
\section*{Tài liệu tham khảo}
\addcontentsline{toc}{section}{Tài liệu tham khảo}
% ══════════════════════════════════════════════════════════════════════════

\begin{thebibliography}{99}

\bibitem{fourier1827}
J. Fourier,
\textit{Solution d'une question particulière du calcul des inégalités},
Nouveau Bulletin des Sciences par la Société Philomathique de Paris, 1827.

\bibitem{motzkin1936}
T. S. Motzkin,
\textit{Beiträge zur Theorie der linearen Ungleichungen},
Dissertation, Universität Basel, 1936.

\bibitem{schrijver1986}
A. Schrijver,
\textit{Theory of Linear and Integer Programming},
John Wiley \& Sons, 1986.

\bibitem{dantzig1963}
G. B. Dantzig,
\textit{Linear Programming and Extensions},
Princeton University Press, 1963.

\bibitem{cormen2009}
T. H. Cormen, C. E. Leiserson, R. L. Rivest, C. Stein,
\textit{Introduction to Algorithms}, 3rd ed.,
MIT Press, 2009. Chapter 29: Linear Programming.

\bibitem{sympy}
SymPy Development Team,
\textit{SymPy: Python library for symbolic mathematics},
\url{https://www.sympy.org}, 2024.

\bibitem{flask}
A. Ronacher et al.,
\textit{Flask: A lightweight WSGI web application framework},
\url{https://flask.palletsprojects.com}, 2024.

\bibitem{chartjs}
Chart.js Contributors,
\textit{Chart.js: Simple yet flexible JavaScript charting},
\url{https://www.chartjs.org}, 2024.

\end{thebibliography}

% ══════════════════════════════════════════════════════════════════════════
\appendix
\section{Phụ lục: Test cases}
% ══════════════════════════════════════════════════════════════════════════

\subsection{Test case chuẩn}

\begin{table}[H]
\centering
\caption{Các test case kiểm thử}
\label{tab:testcases}
\begin{tabularx}{\textwidth}{lXcc}
\toprule
\textbf{TC} & \textbf{Bài toán} & \textbf{Kết quả mong đợi} & \textbf{Trạng thái} \\
\midrule
TC1 & $\max 3x_1+2x_2$, $x_1+x_2\leq4$, $2x_1+x_2\leq6$, $x_i\geq0$
    & $x_1=2, x_2=2, z=10$ & ✓ \\
TC2 & $\max 2x_1+3x_2+x_3$, $x_1+x_2+x_3\leq6$, $x_2+x_3\leq4$, $x_1+x_2\leq4$, $x_i\geq0$
    & $x_1=2, x_2=2, x_3=2, z=12$ & ✓ \\
TC3 & $\min -x_1-2x_2$, $2x_1+x_2\leq6$, $x_1+x_2\geq2$, $-x_1+x_2\geq3$, $x_i\geq0$
    & $x_1=0, x_2=6, z=-12$ & ✓ \\
TC4 & Bài toán vô nghiệm: $x_1\leq-1$, $x_1\geq0$
    & \texttt{feasible=false} & ✓ \\
TC5 & Hệ số phân số: $\frac{1}{3}x_1+\frac{2}{3}x_2\leq1$
    & Kết quả exact & ✓ \\
\bottomrule
\end{tabularx}
\end{table}

\subsection{Cấu trúc JSON response mẫu}

\begin{lstlisting}[caption={JSON response cho TC1}, label={lst:json_response}]
{
  "feasible": true,
  "solution": {"x1": "2", "x2": "2"},
  "solution_numeric": {"x1": 2.0, "x2": 2.0},
  "z": "10",
  "z_numeric": 10.0,
  "obj_type": "max",
  "method": "algebraic",
  "steps_constraints": [
    {
      "buoc": 0,
      "tieu_de": "He ban dau (sau chuan hoa + bien z)",
      "constraints": ["x1 + x2 <= 4", "2x1 + x2 <= 6",
                      "-x1 <= 0", "-x2 <= 0", "-3x1 - 2x2 + z <= 0"],
      "bi_loai": [], "moi_tao": []
    },
    {
      "buoc": 1,
      "tieu_de": "Sau khi loai x1",
      "constraints": [...],
      "bi_loai": [...], "moi_tao": [...]
    }
  ],
  "reasoning": {
    "steps": [
      {"so_buoc": 1, "loai": "INFO", "tieu_de": "Bat dau giai", ...},
      {"so_buoc": 2, "loai": "TINH_TOAN", "cong_thuc": "...", "ket_qua": "..."},
      ...
    ],
    "tong_buoc": 42
  },
  "vertices": [
    {"point": [0,0], "label": "A(0,0)", "z": 0, "is_optimal": false},
    {"point": [2,2], "label": "C(2,2)", "z": 10, "is_optimal": true},
    ...
  ]
}
\end{lstlisting}

\end{document}
